% Week 5 Session B Lab 2 — SCS-CN runoff modeling
% Build: pdflatex Week5_SessionB_Lab2_Report.tex from Lab-reports/ (twice)
% Screenshots:
%   week5_session_b_lab2_terminal1.png  (pytest)
%   week5_session_b_lab2_terminal2.png  (sensitivity_analysis.py)
%   week5_session_b_lab2_terminal3.png  (validate_reference.py)
%   week5_session_b_lab2_sensitivity.png (plot)
% Code: week5_session_b_lab2_files/
%
\documentclass[11pt,a4paper]{article}

\usepackage[margin=2.5cm]{geometry}
\usepackage{graphicx}
\newcommand{\termfig}[1]{%
  \includegraphics[width=\linewidth,height=0.42\textheight,keepaspectratio]{#1}%
}
\newcommand{\plotfig}[1]{%
  \includegraphics[width=\linewidth,height=0.55\textheight,keepaspectratio]{#1}%
}
\usepackage{booktabs}
\usepackage{xurl}
\usepackage{hyperref}
\usepackage{parskip}
\usepackage{enumitem}
\usepackage{xcolor}
\usepackage{fvextra}
\usepackage[most]{tcolorbox}
\usepackage{subcaption}

\newcommand{\studentname}{Mahmudul Hasan}
\newcommand{\studentid}{4125999049}
\newcommand{\reportdate}{May 21, 2026}

\makeatletter
\g@addto@macro\UrlBreaks{\do\ }
\makeatother

\newcommand{\LabCodeRoot}{week5_session_b_lab2_files}

\newcommand{\filebox}[2]{%
\begin{tcolorbox}[
  enhanced,
  colback=gray!6,
  colframe=black!55,
  title={#1},
  fonttitle=\bfseries\ttfamily\footnotesize,
  arc=1.5mm,
  boxrule=0.7pt,
  breakable,
  width=\linewidth,
  before skip=0.8em,
  after skip=0.8em,
]
\VerbatimInput[fontsize=\footnotesize,breaklines=true,breakanywhere=true]{\LabCodeRoot/#2}
\end{tcolorbox}%
}

\title{Lab Report: Week 5 Session B Lab 2\\SCS-CN Runoff Modeling}
\author{\studentname \quad (Student ID: \studentid)}
\date{\reportdate}

\begin{document}
\maketitle

\begin{abstract}
This report documents Week~5 Session~B Lab~2: implementation of the \textbf{SCS-CN} direct runoff method in Python, parameter sensitivity over curve numbers CN$\in\{60,70,80,90,95\}$ and rainfall $P\in[0,100]$~mm, and domain validation ($Q\le P$, $Q=0$ when $P\le I_a$, higher CN $\Rightarrow$ more runoff). The function \texttt{calculate\_runoff(P, CN)} reuses the Week~3 Session~B TDD implementation in \texttt{week5\_session\_b\_lab2/}. All \textbf{7 pytest} cases passed; \texttt{sensitivity\_analysis.py} saved \texttt{figures/runoff\_sensitivity.png} and passed physical checks; reference cases $Q(50,80)\approx 13.80$~mm and $Q(80,85)\approx 43.55$~mm match course expectations. Evidence: three terminal screenshots and the sensitivity figure; source in Appendix~\ref{app:runoff}--\ref{app:promptlog}.
\end{abstract}

\section{Introduction}
The Soil Conservation Service Curve Number (SCS-CN) method estimates rainfall excess (direct runoff) from total storm depth and a land-use--dependent curve number. This lab translates the metric formulas into tested Python code, explores sensitivity of $Q$ to CN, and verifies hydrologic plausibility---skills that complement Week~5 Session~A Lab~1 (live rainfall from an API) with \textbf{process-based} runoff depth in millimetres.

\section{Objectives}
\begin{itemize}[noitemsep]
  \item Implement \texttt{calculate\_runoff(P, CN)} with boundary conditions and validation.
  \item Run comprehensive pytest edge cases (Exercise~1).
  \item Produce sensitivity plots and a pandas DataFrame (Exercise~2).
  \item Validate $Q\le P$, zero runoff below $I_a$, and CN ordering (Exercise~3).
  \item Document work in \texttt{prompt\_log.md} (Exercise~4).
\end{itemize}

\section{Methods}

\subsection{SCS-CN formulation (Exercise 1)}
All depths in \textbf{millimetres (mm)}:
\begin{align*}
S &= \frac{25400}{\mathrm{CN}} - 254, \qquad I_a = 0.2\,S, \\
Q &= \begin{cases}
0 & \text{if } P \le I_a, \\
\dfrac{(P-I_a)^2}{P - I_a + S} & \text{otherwise},
\end{cases}
\end{align*}
with $Q$ clipped to $[0,P]$. CN is validated in $[1,100]$ to avoid division by zero on $S$. Implementation: \texttt{src/runoff.py} (developed in Week~3 Session~B via TDD; reused here).

\subsection{Tests (Exercise 1)}
\texttt{tests/test\_runoff.py} --- seven cases: normal ($P{=}50$, CN${=}80$), $P{=}0$, $P<I_a$, CN${=}100$ (impervious), negative $P$, invalid CN, and grid check $Q\le P$.

\subsection{Sensitivity analysis (Exercise 2)}
\texttt{sensitivity\_analysis.py} builds a wide DataFrame (\texttt{rainfall\_mm}, \texttt{Q\_CN60}, \ldots), saves \texttt{figures/sensitivity\_data.csv}, plots rainfall vs.\ $Q$ for each CN with a dashed $Q=P$ upper bound, and runs embedded domain checks.

\subsection{Reference validation (Exercise 3)}
\texttt{validate\_reference.py} spot-checks course reference points; \texttt{sensitivity\_analysis.domain\_checks()} asserts $Q\le P$ and monotonic increase of $Q$ with CN at $P{=}50$~mm.

\section{Results}

\subsection{File inventory}
\begin{table}[htbp]
  \centering
  \caption{Submitted files in \texttt{week5\_session\_b\_lab2/}.}
  \label{tab:files}
  \small
  \begin{tabular}{@{}ll@{}}
    \toprule
    File & Role \\
    \midrule
    \texttt{src/runoff.py} & SCS-CN \texttt{calculate\_runoff} \\
    \texttt{tests/test\_runoff.py} & pytest suite (7 tests) \\
    \texttt{sensitivity\_analysis.py} & DataFrame, plot, domain checks \\
    \texttt{validate\_reference.py} & Hand-check reference cases \\
    \texttt{figures/runoff\_sensitivity.png} & Sensitivity figure \\
    \texttt{prompt\_log.md}, \texttt{prompt\_*.txt} & Documentation \\
    \bottomrule
  \end{tabular}
\end{table}

\subsection{Terminal verification}
Figures~\ref{fig:pytest}--\ref{fig:validate} show the three commands required for submission evidence.

\begin{figure}[htbp]
  \centering
  \begin{subfigure}[t]{\linewidth}
    \centering
    \termfig{week5_session_b_lab2_terminal1.png}
    \caption{\texttt{python3 -m pytest tests/ -v} --- 7 passed.}
    \label{fig:pytest}
  \end{subfigure}\\[0.6em]
  \begin{subfigure}[t]{\linewidth}
    \centering
    \termfig{week5_session_b_lab2_terminal2.png}
    \caption{\texttt{python3 sensitivity\_analysis.py} --- plot saved; checks PASS.}
    \label{fig:sensitivity-term}
  \end{subfigure}\\[0.6em]
  \begin{subfigure}[t]{\linewidth}
    \centering
    \termfig{week5_session_b_lab2_terminal3.png}
    \caption{\texttt{python3 validate\_reference.py} --- four reference cases OK.}
    \label{fig:validate}
  \end{subfigure}
  \caption{Terminal evidence for Exercises~1--3.}
  \label{fig:terminal-all}
\end{figure}

Key numeric outputs from \texttt{sensitivity\_analysis.py} at $P{=}50$~mm:
\begin{verbatim}
Q values: {60: 1.4, 70: 5.81, 80: 13.8, 90: 27.11, 95: 36.9}
Check 3: P=80 CN=85 -> Q=43.55 mm (course reference ~43.6)
\end{verbatim}

\subsection{Sensitivity figure (Exercise 2)}
Figure~\ref{fig:plot} shows runoff depth vs.\ rainfall for CN = 60, 70, 80, 90, 95. All curves lie at or below the $Q=P$ line; higher CN yields greater runoff for the same $P$, consistent with more impervious catchments.

\begin{figure}[htbp]
  \centering
  \plotfig{week5_session_b_lab2_sensitivity.png}
  \caption{SCS-CN sensitivity: runoff $Q$ (mm) vs.\ rainfall $P$ (mm) by curve number.}
  \label{fig:plot}
\end{figure}

\subsection{Validation checklist}
\begin{table}[htbp]
  \centering
  \caption{Exercise 3--4 validation summary.}
  \label{tab:validation}
  \begin{tabular}{@{}ll@{}}
    \toprule
    Check & Outcome \\
    \midrule
    pytest 7/7 & PASS (0.01~s) \\
    $Q \le P$ for all $P$, CN & PASS \\
    $Q=0$ when $P < I_a$ (e.g.\ $P{=}10$, CN${=}80$) & PASS \\
    Higher CN $\Rightarrow$ higher $Q$ at $P{=}50$~mm & PASS \\
    $Q(80,85) \approx 43.55$~mm & Matches course reference \\
    $Q(50,100)=50$~mm (impervious) & PASS \\
    \bottomrule
  \end{tabular}
\end{table}

\section{Analysis and validation}

\textbf{Relation to Week~3 Session~B:} The same \texttt{calculate\_runoff} logic was first built under TDD (Red--Green--Refactor). Week~5 Lab~2 adds \textbf{sensitivity visualization} and explicit domain scripts rather than new formula changes.

\textbf{YELLOW/RED rainfall alerts (Week~5 Lab~1):} Lab~1 used OpenWeather mm/h thresholds; Lab~2 uses cumulative storm depth $P$ in mm for the SCS-CN equation---different physical quantities, both valid in a Smart Water course.

\textbf{CN bounds:} Minimum CN${=}1$ prevents singular $S$; CN${=}100$ gives $S{=}0$, $I_a{=}0$, so $Q{=}P$ for storms exceeding zero depth---a useful impervious limit test.

\textbf{Optional OpenCode:} Prompts \texttt{prompt\_implement.txt} and \texttt{prompt\_sensitivity.txt} are available; baseline code and tests were already complete from Week~3, so OpenCode was not required for a passing run.

\section{Discussion}
\begin{enumerate}[noitemsep]
  \item A single pytest screenshot plus sensitivity and validate screenshots give clear, gradable evidence without rerunning commands.
  \item The sensitivity plot makes the CN effect visible; the $Q=P$ reference line reinforces the physical cap enforced in code.
  \item Grid test \texttt{test\_runoff\_never\_exceeds\_rainfall} catches formula slips that hand checks might miss.
  \item Live YELLOW/RED wet-weather cases are not needed here; dry API cities in Lab~1 and formula-based Lab~2 cover complementary validation styles.
\end{enumerate}

\section{Problems encountered}
None blocking: pytest and scripts ran on Python~3.8.10 with pytest~4.6.9. All sampled cities in Lab~1 were dry (GREEN); Lab~2 does not depend on weather API data.

\section{Lessons learned}
Keeping $P$, $S$, $I_a$, and $Q$ in mm avoids unit bugs. Tests written before or alongside implementation (Week~3 TDD) make sensitivity scripts safe to extend. Document reference pairs such as $(P,CN)=(80,85)$ in \texttt{prompt\_log.md} so graders can compare to published examples quickly.

\section{Conclusion}
Week~5 Session~B Lab~2 objectives were met: SCS-CN implementation with tests, sensitivity figure for CN 60--95, and domain validation including course reference runoff $\approx 43.55$~mm at $P{=}80$~mm, CN${=}85$. Submitted: \texttt{week5\_session\_b\_lab2/}, \texttt{prompt\_log.md}, terminal screenshots (Figures~\ref{fig:pytest}--\ref{fig:validate}), and sensitivity plot (Figure~\ref{fig:plot}).

\appendix

\section{Runoff module: \texttt{runoff.py}}
\label{app:runoff}
\filebox{\texttt{src/runoff.py} (full file)}{runoff.py}

\section{Tests: \texttt{test\_runoff.py}}
\label{app:tests}
\filebox{\texttt{tests/test\_runoff.py} (full file)}{test_runoff.py}

\section{Sensitivity script: \texttt{sensitivity\_analysis.py}}
\label{app:sensitivity}
\filebox{\texttt{sensitivity\_analysis.py} (full file)}{sensitivity_analysis.py}

\section{Reference checks: \texttt{validate\_reference.py}}
\label{app:validate}
\filebox{\texttt{validate\_reference.py} (full file)}{validate_reference.py}

\section{Submitted prompt log: \texttt{prompt\_log.md}}
\label{app:promptlog}
\filebox{\texttt{prompt\_log.md} (full file)}{prompt_log.md}

\end{document}
